\documentclass{article}

\usepackage[english]{babel}
\usepackage[letterpaper,top=2cm,bottom=2cm,left=3cm,right=3cm,marginparwidth=1.75cm]{geometry}
\usepackage{amsmath}
\usepackage{amssymb}
\usepackage{mathtools}
\usepackage{amsthm}
\usepackage[colorlinks=true, allcolors=blue]{hyperref}
\usepackage[noabbrev]{cleveref}

\newcommand{\Fp}{\mathbb{F}_p}
\newcommand{\Fq}{\mathbb{F}_q}
\newcommand{\extdeg}{d}
\newcommand{\pow}{\mathsf{pow}}
\newcommand{\fold}{\mathsf{fold}}

\newtheorem{lemma}{Lemma}
\newtheorem{theorem}{Theorem}
\newtheorem{proposition}{Proposition}
\newtheorem{corollary}{Corollary}
\theoremstyle{definition}
\newtheorem{definition}{Definition}
\newtheorem{remark}{Remark}
\newtheorem{condition}{Condition}
\crefname{condition}{condition}{conditions}
\Crefname{condition}{Condition}{Conditions}

\title{LeanVM: Zero-Knowledge}
\author{draft (Work In Progress, not implemented)}
\date{}

\begin{document}
\maketitle

\section{Introduction}

This document is a companion to \texttt{minimal\_zkVM.tex}. It makes the LeanVM proving system zero-knowledge. Notation, the WHIR protocol, the stacking, and the constants are as defined there.

\paragraph{Goal.} Statistical honest-verifier zero-knowledge at $\lambda$ bits: a simulator, given only the statement, must produce transcripts at statistical distance at most $2^{-\lambda}$ from those of the honest prover. We treat the \emph{interactive} protocol. TODO make \emph{BCS transform}~\cite{BCS} (compilation it into a non-interactive argument via Merkle commitments and Fiat--Shamir) ZK.

\paragraph{Result.} The WHIR opening leaks: its sumcheck messages, query answers, and final polynomial are linear in the committed data. We mask the commitment (\Cref{sec:zk-whir}) at a prover cost slightly above $2\times$, with no new prover-side check and no change for the verifier beyond the requirement $s_0 \geq 2$ on the number of out-of-domain samples at commitment. \Cref{sec:zk-proof} proves that the masked opening is zero-knowledge: on an explicit \emph{good set} of challenge tuples, determined by the challenges and the statement alone, the simulation is \emph{perfect} --- the simulator is the honest prover run on substitute data computed from the statement (\Cref{thm:simulator}) --- and the complement of the good set has probability
\[
\Pr[\text{degenerate challenges}] \;\leq\; \big(s_0 + s_1 + 1\big)\, \frac{p}{q} \;+\; O\big(2^{k_0+7}\big)\frac{1}{q} \;\approx\; 2^{-122} \quad \text{for KoalaBear}
\]
(\Cref{cor:restart,cor:statzk}). The claimed value $c$, input of the opening, must be made zero-knowledge separately; this is the subject of \Cref{sec:zk-logup-gkr,sec:zk-air-sumcheck,sec:zk-column-evals}, partly open.

\section{The masked WHIR}\label{sec:zk-whir}

Three changes to the protocol.

\paragraph{1. Folding order.} We use the folding ordering $(X_1, X_n, X_{n-1}, \dots, X_2)$: round 0 folds $X_1$ first, then $X_n, X_{n-1}, \dots, X_{n-k_0+2}$; the later rounds continue downward; the $m_M$ variables surviving to the final polynomial are $(X_2, \dots, X_{m_M + 1})$.

\paragraph{2. Layout.} Two uniformly random \emph{masks} are committed next to the data:
\[
P \;=\; \underbrace{\,\text{stacking} \;\Vert\; 0 \cdots 0 \;\Vert\; R_{\mathsf{in}}\,}_{P(0,\, \cdot\,) \text{ : the data half}} \;\Big\Vert\; \underbrace{\,R_{\mathsf{out}}\,}_{P(1,\, \cdot\,)}
\]
\begin{itemize}
    \item The stacking part (containing the secret data) is unchanged: we simply commit to a bigger polynomial, to fit $R_{\mathsf{in}}$ and $R_{\mathsf{out}}$.
    \item The mask $R_{\mathsf{out}}$ fills the second half with i.i.d. uniform $\Fp$ cells.
    \item The mask $R_{\mathsf{in}}$ is a uniformly random segment occupying the \emph{last} $2^{k_0 - 1 + a}$ cells of the first half of the polynomial, with $2^a \geq t_0 + \extdeg \cdot (s_0 + s_1)$ (a negligible size in practice).
\end{itemize}

\paragraph{3. Out-of-domain count.} We require $s_0 \geq 2$ out-of-domain points at commitment. This is a zero-knowledge requirement, not necessarily a soundness one: with $s_0 = 1$, the round-0 sumcheck messages can never be masked, whatever the size of the masks (\Cref{rem:identities}).

There is no further change. In particular, the prover performs no check, and soundness is unaffected: the masks occupy cells where every claim weight vanishes.

\section{Zero-knowledge of the masked WHIR}\label{sec:zk-proof}

\subsection{Setting}\label{sec:setting}

\paragraph{Cells and classes.} The committed vector $P \in \Fp^{2^n}$ is indexed by $u \in \{0,1\}^n$, the coordinates listed in folding order: first the $k_0$ \emph{round-0 coordinates}, then the $n - k_0$ \emph{position coordinates}. For $s \in \{0,1\}^{k_0}$, the \emph{class} $u_s \in \Fp^{2^{n - k_0}}$ is the vector of cells with round-0 coordinates $s$, indexed by position. The layout of \Cref{sec:zk-whir} reads as follows on the classes: a class with $s_1 = 1$ lies entirely inside $R_{\mathsf{out}}$; a class with $s_1 = 0$ contains the data and the zero padding, and ends with $2^a$ cells of $R_{\mathsf{in}}$, at its last $2^a$ positions. By locality of the fold, the $2^{k_0}$ values of $f_0$ opened at a STIR query $x$ carry the same information as the per-class evaluations $\big(\hat{u}_s(\pow(x, n - k_0))\big)_s$; we work with the latter throughout. The round-0 fold acts classwise: at every position $c$,
\[
\hat{f}_1[c] \;=\; \sum_{s \in \{0,1\}^{k_0}} \hat{eq}(\vec{\alpha}_0, s)\, u_s[c] ,
\]
where $\vec{\alpha}_0 = (\alpha_1, \dots, \alpha_{k_0})$ are the round-0 sumcheck challenges, in folding order.

\paragraph{Points and nodes.} For a point $z \in \Fq$, in folding order, $\pow(z, n) = (\pow(z, k_0), \nu)$, where $\nu := \pow(z^{2^{k_0}}, n - k_0)$ is the \emph{node} of $z$.

\paragraph{Weight and messages.} We work with a single point until \Cref{sec:twopoint}, which batches two, and with the one-point weight $\hat{W}_0 = \hat{eq}(\pow(z, n), \cdot) + \gamma\, \hat{w}$. The input weight $\hat{w}$ vanishes on every mask cell. The round-0 sumcheck messages are
\[
\hat{h}_\ell(X) \;=\; \sum_{b \in \{0,1\}^{n - \ell}} \hat{f}_0(\alpha_1, \dots, \alpha_{\ell - 1}, X, b)\; \hat{W}_0(\alpha_1, \dots, \alpha_{\ell - 1}, X, b) , \qquad \ell = 1, \dots, k_0 ,
\]
each of degree $2$ in $X$, with two fresh coefficients (the third is fixed by $\hat{h}_\ell(0) + \hat{h}_\ell(1) = \sigma$). The \emph{round-0 view} consists of: the per-class evaluations at the $t_0$ distinct queried points; the out-of-domain answers; the $2k_0$ fresh sumcheck coefficients; and the $s_1$ out-of-domain answers on $\hat{f}_1$.

\subsection{The master polynomial and the node matrix}\label{sec:master}

Summing the position coordinates first in the definition of $\hat{h}_\ell$ leaves a partial boolean sum of a single polynomial in the $k_0$ round-0 variables:
\[
\hat{H}(Y) \;:=\; \sum_{c \in \{0,1\}^{n - k_0}} \hat{f}_0(Y, c)\, \hat{W}_0(Y, c) , \qquad \hat{h}_\ell(X) \;=\; \sum_{b \in \{0,1\}^{k_0 - \ell}} \hat{H}(\alpha_1, \dots, \alpha_{\ell - 1}, X, b) .
\]
In the out-of-domain term of $\hat{H}$, the weight factorizes coordinate by coordinate, $\hat{eq}(\pow(z, n), (Y, c)) = \hat{eq}(\pow(z, k_0), Y)\, \hat{eq}(\nu, c)$, and summing $\hat{eq}(\nu, \cdot)$ against the multilinear $\hat{f}_0$ collapses the position coordinates onto the node:
\begin{equation}\label{eq:master}
\hat{H}(Y) \;=\; \hat{eq}(\pow(z, k_0), Y)\, \hat{f}_0(Y, \nu) \;+\; \gamma \sum_{c} \hat{f}_0(Y, c)\, \hat{w}(Y, c) , \qquad \hat{f}_0(Y, \nu) \;=\; \sum_s \hat{eq}(Y, s)\, \hat{u}_s(\nu) .
\end{equation}
The node restriction $\hat{f}_0(\cdot, \nu)$ is multilinear in the $k_0$ round-0 variables, with hypercube values the \emph{node values} $\hat{u}_s(\nu)$; the out-of-domain answer is $y = \hat{f}_0(\pow(z, n))$, and the node evaluation of the folded polynomial is $\hat{f}_1(\nu) = \hat{f}_0(\vec{\alpha}_0, \nu)$.

\paragraph{The masked channel.} Substituting \cref{eq:master} into $\hat{h}_\ell$ and resumming the trailing boolean coordinates by multilinearity, the out-of-domain contribution to message $\ell$ is
\begin{equation}\label{eq:channel}
P_\ell\; \hat{eq}(z^{2^{\ell-1}}, X)\; g_\ell(X) , \qquad P_\ell := \prod_{i < \ell} \hat{eq}(\alpha_i, z^{2^{i-1}}) ,
\end{equation}
with $g_\ell(X) := \hat{f}_0(\alpha_1, \dots, \alpha_{\ell - 1},\, X,\, z^{2^{\ell}}, \dots, z^{2^{k_0-1}},\, \nu)$, an \emph{affine} polynomial (multilinear, one free variable). Up to the public factors $P_\ell\, \hat{eq}(z^{2^{\ell-1}}, X)$, message $\ell$ shows $g_\ell$, and nothing more. Directly from the definition of $g_\ell$, the $g_\ell$ form a \emph{chain}:
\begin{equation}\label{eq:chain}
g_1(z) \;=\; y , \qquad g_{\ell+1}(z^{2^{\ell}}) \;=\; g_\ell(\alpha_\ell) \quad (\ell < k_0) , \qquad g_{k_0}(\alpha_{k_0}) \;=\; \hat{f}_0(\vec{\alpha}_0, \nu) \;=\; \hat{f}_1(\nu) .
\end{equation}
Each affine $g_\ell$ has two coefficients, one of which is chained to the previous round: the masked content of round 0 is $y$ plus \emph{one fresh dimension per round}, $k_0 + 1$ in total. \Cref{sec:rank} proves that this count is exact.

\paragraph{The node matrix.} Expanding $g_\ell$ in the node values turns \cref{eq:channel} into $\sum_s \Lambda_\ell(X; s)\, \hat{u}_s(\nu)$, with
\begin{equation}\label{eq:lambda}
\Lambda_\ell(X; s) \;=\; \Big[\prod_{i < \ell} \hat{eq}(\alpha_i, s_i)\, \hat{eq}(\alpha_i, z^{2^{i-1}})\Big] \cdot \hat{eq}(X, s_\ell)\, \hat{eq}(X, z^{2^{\ell-1}}) \cdot \prod_{i = \ell + 1}^{k_0} \hat{eq}(s_i, z^{2^{i-1}}) .
\end{equation}
We study the following $2 k_0 + 2$ linear forms in the $2^{k_0}$ node values $\hat{u}_s(\nu)$, the \emph{node matrix}:
\begin{itemize}
    \item the out-of-domain answer $y$: the $\omega$-row, with coefficient $\prod_{i \le k_0} \hat{eq}(z^{2^{i-1}}, s_i)$ at column $s$;
    \item the fresh sumcheck coefficients: rows $(\ell, r)$ for $\ell \le k_0$, $r \in \{1, 2\}$, with coefficient $[X^r]\, \Lambda_\ell(X; s)$ at column $s$;
    \item the node evaluation $\hat{f}_1(\nu)$: the $\eta$-row, with coefficient $\hat{eq}(\vec{\alpha}_0, s)$ at column $s$.
\end{itemize}

\paragraph{The $\hat{w}$ term and the cross components.} Expanding $\hat{f}_0 = \sum_u P[u]\, \hat{eq}(u, \cdot)$ instead, the $\hat{w}$ term contributes to $\hat{h}_\ell$ the per-cell weight
\[
\gamma \cdot \Big[\prod_{i < \ell} \hat{eq}(\alpha_i, u_i)\Big] \cdot \hat{eq}(X, u_\ell) \cdot \hat{w}(\alpha_1, \dots, \alpha_{\ell-1}, X, u_{> \ell})
\]
at the cell $u$. This weight vanishes at every cell of $R_{\mathsf{in}}$ (\Cref{lem:blocks} below). Summing up: as a linear form on the cells, every revealed value $v$ (an out-of-domain answer, or a fresh coefficient of a sumcheck message) is the sum of three parts. The out-of-domain term contributes the \emph{node component} $N_v$, a linear form in the node values $\big(\hat{u}_s(\nu)\big)_s$. The $\hat{w}$ term contributes a form supported on the data cells and on $R_{\mathsf{out}}$; the $R_{\mathsf{out}}$ part is the \emph{cross component} $C_v$, and its contributions are called the \emph{cross-terms}.

\subsection{Slot decomposition}\label{sec:slot}

A linear form in the node values that factorizes per coordinate of $s$ is a tensor $v_1 \otimes \cdots \otimes v_{k_0}$ of \emph{slot vectors} $v_i \in \Fq^2$, where $v_i = (v_i[0], v_i[1])$ collects the two weights of the coordinate $s_i$. Define
\[
c_i := (1 - z^{2^{i-1}},\; z^{2^{i-1}}), \qquad \tilde{a}_i := (1 - \alpha_i,\; \alpha_i), \qquad d := (-1,\; 1) .
\]
Then $\omega = c_1 \otimes \cdots \otimes c_{k_0}$ and $\eta = \tilde{a}_1 \otimes \cdots \otimes \tilde{a}_{k_0}$. For the sumcheck rows, write $\hat{eq}(X, t) = (1 - t) + X(2t - 1)$ and expand the slot-$\ell$ pair $\big((1-X)\hat{eq}(X, z^{2^{\ell-1}}),\; X \hat{eq}(X, z^{2^{\ell-1}})\big)$ of \cref{eq:lambda}:
\[
b^{(1)}_\ell := \big(3z^{2^{\ell-1}} - 2,\; 1 - z^{2^{\ell-1}}\big) \quad (X^1 \text{ coefficients}), \qquad
b^{(2)}_\ell := (2z^{2^{\ell-1}} - 1) \cdot (-1, 1) \quad (X^2 \text{ coefficients}),
\]
so that
\[
\text{row}(\ell, r) \;=\; P_\ell \cdot \tilde{a}_1 \otimes \cdots \otimes \tilde{a}_{\ell - 1} \otimes b^{(r)}_\ell \otimes c_{\ell + 1} \otimes \cdots \otimes c_{k_0} .
\]
Three identities, each a direct computation:
\begin{align}
\det(c_i, d) &= 1 \qquad \text{(so $\{c_i, d\}$ is always a basis of the slot plane)}, \label{eq:cd}\\
\tilde{a}_i &= c_i + (\alpha_i - z^{2^{i-1}})\, d , \label{eq:acd}\\
b^{(1)}_\ell = (2z^{2^{\ell-1}} - 1)\, c_\ell + (1 - 2z^{2^{\ell}})\, d , &\qquad b^{(2)}_\ell = (2z^{2^{\ell-1}} - 1)\, d , \qquad \det\big(b^{(1)}_\ell, b^{(2)}_\ell\big) = (2z^{2^{\ell-1}} - 1)^2 . \label{eq:bcd}
\end{align}
Define, for $\ell = 1, \dots, k_0 + 1$ and $m = 1, \dots, k_0$:
\[
\rho_\ell := \tilde{a}_1 \otimes \cdots \otimes \tilde{a}_{\ell-1} \otimes c_\ell \otimes \cdots \otimes c_{k_0}, \qquad
\tau_m := \tilde{a}_1 \otimes \cdots \otimes \tilde{a}_{m-1} \otimes d \otimes c_{m+1} \otimes \cdots \otimes c_{k_0},
\]
with the conventions $\rho_1 = \omega$ and $\rho_{k_0 + 1} = \eta$. By \cref{eq:bcd},
\begin{equation}\label{eq:rows}
\text{row}(\ell, 1) = P_\ell \big[(2z^{2^{\ell-1}} - 1)\rho_\ell + (1 - 2z^{2^{\ell}})\tau_\ell\big], \qquad \text{row}(\ell, 2) = P_\ell (2z^{2^{\ell-1}} - 1)\, \tau_\ell .
\end{equation}
By \cref{eq:acd}, replacing $\tilde{a}_\ell$ by $c_\ell + (\alpha_\ell - z^{2^{\ell-1}})d$ in slot $\ell$ of $\rho_{\ell + 1}$ gives the \emph{telescoping identity}
\begin{equation}\label{eq:telescope}
\rho_{\ell + 1} \;=\; \rho_\ell + (\alpha_\ell - z^{2^{\ell-1}})\, \tau_\ell , \qquad \ell = 1, \dots, k_0 .
\end{equation}

\subsection{Rank and kernel at one point}\label{sec:rank}

This subsection is the per-point engine of the two-point analysis: \Cref{thm:twopoint} consumes the rank and the explicit kernel basis of \Cref{prop:rank} (per point), \Cref{cor:twopointprob} consumes \Cref{cor:rankprob}, and the twist-rigidity argument (\Cref{cond:twist}) uses $\eta \in \mathrm{span}\{\omega, \tau_1, \dots, \tau_{k_0}\}$.

\begin{lemma}\label{lem:tau}
The $k_0 + 1$ tensors $\omega, \tau_1, \dots, \tau_{k_0}$ are linearly independent, for all values of the variables.
\end{lemma}

\begin{proof}
Induction on $k_0$. For $k_0 = 1$: $\omega = c_1$ and $\tau_1 = d$ are independent by \cref{eq:cd}. For $k_0 \ge 2$, let $\omega', \tau'_1, \dots, \tau'_{k_0 - 1}$ be the corresponding tensors on the slots $2, \dots, k_0$. Then
\[
\omega = c_1 \otimes \omega', \qquad \tau_1 = d \otimes \omega', \qquad \tau_m = \tilde{a}_1 \otimes \tau'_{m - 1} \ (m \ge 2) .
\]
Suppose $\lambda_\omega \omega + \sum_m \lambda_m \tau_m = 0$. Using \cref{eq:acd} to express $\tilde{a}_1$ in the basis $\{c_1, d\}$ and collecting the two basis components, with $Y := \sum_{m \ge 2} \lambda_m \tau'_{m-1}$:
\[
\lambda_\omega\, \omega' + Y = 0 \qquad \text{and} \qquad \lambda_1\, \omega' + (\alpha_1 - z)\, Y = 0 .
\]
By the induction hypothesis, $\omega', \tau'_1, \dots, \tau'_{k_0-1}$ are independent; since $Y$ has no $\omega'$ component, the first equation forces $\lambda_\omega = 0$ and $\lambda_m = 0$ for $m \ge 2$, and then the second forces $\lambda_1 = 0$.
\end{proof}

\begin{proposition}\label{prop:rank}
Assume $z^{2^{\ell-1}} \neq 1/2$ for $\ell \le k_0$ and $\hat{eq}(\alpha_i, z^{2^{i-1}}) \neq 0$ for $i \le k_0 - 1$. Then the node matrix has rank exactly $k_0 + 1$: its row space equals $\mathrm{span}\{\omega, \tau_1, \dots, \tau_{k_0}\}$, and a basis of it is given by the rows $\omega, \mathrm{row}(1,2), \dots, \mathrm{row}(k_0, 2)$. The $k_0 + 1$ remaining rows satisfy the explicit dependencies obtained from \cref{eq:rows,eq:telescope}: for each $\ell$, $\mathrm{row}(\ell, 1)$ lies in the span of $\omega$ and $\mathrm{row}(1,2), \dots, \mathrm{row}(\ell, 2)$, and $\eta = \omega + \sum_{\ell \le k_0} (\alpha_\ell - z^{2^{\ell-1}})\, \tau_\ell$.
\end{proposition}

\begin{proof}
By \cref{eq:rows}, $\mathrm{row}(\ell, 2) = P_\ell (2z^{2^{\ell-1}} - 1) \tau_\ell$ with $P_\ell(2z^{2^{\ell-1}} - 1) \neq 0$ under the assumptions, so $\{\omega, \mathrm{row}(\ell,2)_\ell\}$ spans $\mathrm{span}\{\omega, \tau_1, \dots, \tau_{k_0}\}$, which has dimension $k_0 + 1$ by \Cref{lem:tau}. Iterating \cref{eq:telescope} from $\rho_1 = \omega$ gives $\rho_\ell = \omega + \sum_{i < \ell}(\alpha_i - z^{2^{i-1}})\tau_i$ for every $\ell \le k_0 + 1$; in particular every $\rho_\ell$, hence every $\mathrm{row}(\ell, 1)$ by \cref{eq:rows}, and $\eta = \rho_{k_0+1}$, lie in that same span.
\end{proof}

\begin{corollary}[failure probability]\label{cor:rankprob}
For uniform $z \in \Fq$ and uniform challenges, the assumptions of \Cref{prop:rank} fail with probability at most $2^{k_0}/q + (k_0 - 1)/q$. Indeed $z^{2^{\ell - 1}} = 1/2$ has at most $2^{\ell-1}$ solutions, for a total of at most $2^{k_0} - 1$; and for fixed $z$, the equation $\hat{eq}(\alpha_i, z^{2^{i-1}}) = 0$ is affine non-degenerate in $\alpha_i$, with exactly one solution.
\end{corollary}

\begin{remark}[the dependencies are the chain]\label{rem:identities}
The dependencies of \Cref{prop:rank} are \cref{eq:chain} in matrix form, and hold for \emph{every} choice of challenges. The telescoping identity \cref{eq:telescope} is the link $g_{\ell+1}(z^{2^{\ell}}) = g_\ell(\alpha_\ell)$; the case $\ell = 1$, $\rho_1 = \omega$, is the link $g_1(z) = y$: the first sumcheck message of a batched evaluation claim \emph{contains} that claim. The identity $\eta = \rho_{k_0+1}$ is the exit $g_{k_0}(\alpha_{k_0}) = \hat{f}_1(\nu)$, equivalently the terminal sumcheck claim $\sum_b \hat{f}_1(b)\, \hat{W}_0(\vec{\alpha}_0, b) = \hat{h}_{k_0}(\alpha_{k_0})$ restricted to node components. Consequently, with a single out-of-domain point, no choice of challenges makes the $2k_0 + 2$ node rows independent, and no mask size compensates a rank deficit: this is why the design requires $s_0 \geq 2$.
\end{remark}

\subsection{Two points: full rank}\label{sec:twopoint}

Now take $s_0 = 2$ points $z_1, z_2$, with nodes $\nu_j := \pow(z_j^{2^{k_0}}, n - k_0)$, and node values $\hat{u}_s(\nu_j)$, treated as $2 \cdot 2^{k_0}$ independent variables $V_{s,j}$ (columns indexed by pairs $(s, j)$). The batched weight is $\hat{W}_0 = \hat{eq}(\pow(z_1, n), \cdot) + \gamma\, \hat{eq}(\pow(z_2, n), \cdot) + \gamma^2 \hat{w}$: the master polynomial now carries one channel per point, as in \cref{eq:master}, and every message superposes two chains. The transcript values, as linear forms in the $V_{s,j}$, are the \emph{value rows}:
\begin{itemize}
    \item the two out-of-domain answers $\omega_1, \omega_2$: answer $j$ is supported on the columns $(\cdot, j)$, where it equals the one-point row $\omega^{(j)}$ of \Cref{sec:slot} with $z = z_j$;
    \item the $2k_0$ fresh sumcheck coefficients: the value row $(\ell, r)$ has block-$j$ component $\gamma^{j-1}\, \mathrm{row}(\ell, r)^{(j)}$, on both blocks simultaneously.
\end{itemize}
We keep the name \emph{node matrix} for the resulting $(2k_0 + 2) \times (2 \cdot 2^{k_0})$ matrix of value rows.

\begin{theorem}\label{thm:twopoint}
Let $g(x) := \dfrac{1 - 2x^2}{2x - 1}$. Assume $\gamma \neq 0$, the assumptions of \Cref{prop:rank} for each point, and, for each slot $\ell \leq k_0$:
\[
z_1^{2^{\ell-1}} \neq z_2^{2^{\ell-1}} \qquad \text{and} \qquad \big(2z_1^{2^{\ell-1}} - 1\big)\big(2z_2^{2^{\ell-1}} - 1\big) \neq -1 .
\]
Then the $2k_0 + 2$ value rows are linearly independent. Moreover, their span contains the two extension rows $\eta_j$ (coefficient $\hat{eq}(\vec{\alpha}_0, s)$ at column $(s,j)$, supported on block $j$): the node combinations giving the evaluations of $\hat{f}_1$ at the two nodes are determined by the revealed values. These evaluations accordingly join the directions pinned by the view; this is harmless, since pinned values are functions of the view, which the simulator generates itself.
\end{theorem}

\begin{proof}
Suppose $\mu_1 \omega_1 + \mu_2 \omega_2 + \sum_{\ell, r} \nu_{\ell,r}\, (\text{value row } (\ell,r)) = 0$. Restricting to the columns of block $j$:
\[
\mu_j\, \omega^{(j)} + \gamma^{j-1} \sum_{\ell, r} \nu_{\ell,r}\, \mathrm{row}(\ell, r)^{(j)} \;=\; 0 .
\]
Consider the left kernel $K^{(j)}$ of the $(2k_0+1)$-row one-point matrix $\big[\omega^{(j)}; \mathrm{row}(\ell,r)^{(j)}\big]$. By \Cref{prop:rank} its rank is $k_0 + 1$, so $\dim K^{(j)} = k_0$, and \cref{eq:rows,eq:telescope} give an explicit basis $\{k^{(j)}_\ell\}_{\ell \leq k_0}$: writing $\tau^{(j)}_\ell = \mathrm{row}(\ell,2)^{(j)} / (P^{(j)}_\ell (2z_j^{2^{\ell-1}} - 1))$ and substituting the telescoped $\rho^{(j)}_\ell$ into $\mathrm{row}(\ell,1)^{(j)}$,
\[
\begin{aligned}
k^{(j)}_\ell : \quad & 1 \ \text{on } (\ell,1), \qquad -\,g\big(z_j^{2^{\ell-1}}\big) \ \text{on } (\ell,2), \qquad -\,P^{(j)}_\ell \big(2z_j^{2^{\ell-1}} - 1\big) \ \text{on } \omega , \\
& -\,\frac{P^{(j)}_\ell \big(2z_j^{2^{\ell-1}} - 1\big)\big(\alpha_i - z_j^{2^{i-1}}\big)}{P^{(j)}_i \big(2z_j^{2^{i-1}} - 1\big)} \ \text{on } (i,2) \ \text{for } i < \ell .
\end{aligned}
\]
The projection $\pi$ dropping the $\omega$-coordinate is injective on $K^{(j)}$: an element with all $(\ell,r)$-coordinates zero has $\mu\, \omega^{(j)} = 0$, and $\omega^{(j)} \neq 0$. The two block restrictions say $(\mu_1, \nu) \in K^{(1)}$ and $(\mu_2, \gamma \nu) \in K^{(2)}$; since $\gamma \neq 0$ and $\pi(K^{(2)})$ is a subspace, $\nu \in \pi(K^{(1)}) \cap \pi(K^{(2)})$.

Write $\nu = \sum_\ell \lambda_\ell\, \pi(k^{(1)}_\ell) = \sum_\ell \lambda'_\ell\, \pi(k^{(2)}_\ell)$. The $(\ell, 1)$-coordinates give $\lambda_\ell = \lambda'_\ell$ for every $\ell$. The $(i, 2)$-coordinates then give, for each $i$,
\[
\lambda_i \, \big( g(z_1^{2^{i-1}}) - g(z_2^{2^{i-1}}) \big) \;=\; \sum_{\ell > i} \lambda_\ell \, \Big( c^{(2)}_{\ell, i} - c^{(1)}_{\ell, i} \Big) ,
\]
with $c^{(j)}_{\ell,i}$ the $(i,2)$-coefficients above: an upper-triangular homogeneous system with diagonal $d_i := g(z_1^{2^{i-1}}) - g(z_2^{2^{i-1}})$. A direct computation gives
\[
g(x) - g(y) \;=\; \frac{(y - x)\,\big[(2x-1)(2y-1) + 1\big]}{(2x-1)(2y-1)} ,
\]
so the assumptions say exactly $d_i \neq 0$ for all $i$. Back-substitution from $i = k_0$ downward gives $\lambda = 0$, hence $\nu = 0$, hence $\mu_j\, \omega^{(j)} = 0$ and $\mu_j = 0$.

For the extension rows: by \Cref{lem:tau} applied per block, the forms $\{\omega^{(j)}, \tau^{(j)}_1, \dots, \tau^{(j)}_{k_0}\}_{j = 1,2}$ span a space of dimension exactly $2(k_0+1)$, which contains all value rows. The $2k_0 + 2$ independent value rows therefore span it entirely; and $\eta_j = \rho^{(j)}_{k_0+1}$ lies in it by \Cref{prop:rank}.
\end{proof}

\begin{corollary}\label{cor:twopointprob}
For uniform independent $z_1, z_2, \gamma$ and challenges, the assumptions of \Cref{thm:twopoint} fail with probability at most
\[
\frac{1}{q} \;+\; \frac{2\,(2^{k_0} + k_0 - 1)}{q} \;+\; \frac{2\,(2^{k_0} - 1)}{q} \;\leq\; \frac{4 \cdot 2^{k_0} + 2 k_0}{q} .
\]
Indeed: $\gamma = 0$ costs $1/q$; the per-point assumptions cost $(2^{k_0} + k_0 - 1)/q$ each (\Cref{cor:rankprob}); for each slot, $z_1^{2^{\ell-1}} = z_2^{2^{\ell-1}}$ has at most $2^{\ell-1}$ solutions in $z_2$ for fixed $z_1$, and $(2z_1^{2^{\ell-1}} - 1)(2z_2^{2^{\ell-1}} - 1) = -1$ determines $z_2^{2^{\ell-1}}$, leaving at most $2^{\ell - 1}$ preimages; both sum to $2(2^{k_0} - 1)/q$ over the slots.
\end{corollary}

\begin{remark}[numerical verification]
At $k_0 = 7$ with random challenges in the KoalaBear field: the $16$ value rows have rank $16$ (full), and adding the extension rows does not increase the rank; with a single point, the $15$ value rows have rank $8 = k_0 + 1$, matching \Cref{prop:rank}. Both checks were run with the actual powers $z^{2^{i-1}}$ and with independent values in their place.
\end{remark}

\begin{remark}[$s_0 > 2$]\label{rem:s0}
The analysis is written for $s_0 = 2$, the minimal and tightest case: the $2k_0 + 2$ value rows then span the full space $\bigoplus_j \mathrm{span}\{\omega_j, \tau_{1,j}, \dots, \tau_{k_0,j}\}$, of dimension exactly $2(k_0 + 1)$, which is what pins the extension rows (the second part of \Cref{thm:twopoint}) and shapes \Cref{cond:cross2,prop:pinbound}. For $s_0 > 2$ the masking only gains channels. (i) Joint independence of the $s_0 + 2k_0$ value rows holds under the conditions of \Cref{thm:twopoint} applied to any \emph{single pair} of points: restricting a vanishing combination to the two corresponding blocks gives $\nu \in \pi(K^{(1)}) \cap \pi(K^{(2)}) = 0$ as in the proof, and the remaining blocks force $\mu_j\, \omega^{(j)} = 0$; \Cref{prop:uniform} and the budget $2^a \geq t_0 + \extdeg(s_0 + s_1)$ are already stated for general $s_0$. (ii) The value rows now span a \emph{proper} subspace, of codimension $(s_0 - 2)\, k_0$: the extension rows are no longer automatically in their span, and a generic dimension count suggests that the only pinned combination $\sum_j \theta_j \eta_j$ is the terminal one, which the terminal protocol direction accounts for; \Cref{cond:cross2,prop:pinbound} would simplify accordingly, at the cost of a genericity lemma quantifying that count. We do not pursue this: $s_0 = 2$ suffices for zero-knowledge, and extra points only add prover and proof cost.
\end{remark}

\subsection{Cross-free blocks and joint uniformity}\label{sec:uniform}

\begin{definition}[blocks]\label{def:blocks}
The \emph{block} of a class $s$ is its last $2^a$ cells: cells of $R_{\mathsf{in}}$ if $s_1 = 0$, of $R_{\mathsf{out}}$ if $s_1 = 1$. All blocks consist of i.i.d.\ uniform cells, and the input weight $\hat{w}$ vanishes at every cell whose position lies in the last $2^a$ range, on every class.
\end{definition}

Read as a univariate polynomial of degree $< 2^a$ (the \emph{block univariate} of the class), a block contributes $x^{2^{n - k_0} - 2^a}$ times its evaluation at $x$ to the class evaluation $\hat{u}_s(\pow(x, n - k_0))$: above the block, every coordinate of the evaluation point carries the bit $1$, contributing $x^{2^a} x^{2^{a+1}} \cdots x^{2^{n - k_0 - 1}} = x^{2^{n - k_0} - 2^a}$, a prefactor that never vanishes. (This is why the blocks sit at the \emph{end} of each class: a coordinate carrying the bit $0$ would contribute a factor $(1 - x^{2^i})$ instead, which vanishes on small-order $x$, a constant fraction of the queries.)

\begin{lemma}[cross-free blocks]\label{lem:blocks}
For every round $\ell$ and every cell $u$ of a block, the cross-term weight vanishes identically: $\hat{w}(\alpha_1, \dots, \alpha_{\ell-1}, X, u_{>\ell}) = 0$.
\end{lemma}

\begin{proof}
The partial evaluation is the multilinear extension, in the folded coordinates, of the values of $\hat{w}$ at the boolean completions of $u$. A completion changes only round-0 coordinates, hence preserves the position of $u$ inside its class, which lies in the last $2^a$ range: $\hat{w}$ vanishes at every completion, and the multilinear extension of zeros is zero.
\end{proof}

\begin{lemma}[Chinese remaindering]\label{lem:crt}
Let $x_1, \dots, x_{t} \in \Fp \setminus \{0\}$ be distinct and $\nu_1, \dots, \nu_e \in \Fq \setminus \Fp$ be pairwise distinct and pairwise non-conjugate under the Frobenius $y \mapsto y^p$, with $t + \extdeg\, e \leq 2^a$. For a uniform $g \in \Fp^{2^a}$, viewed as a univariate polynomial of degree $< 2^a$, the tuple
\[
\big(g(x_1), \dots, g(x_{t}),\; g(\nu_1), \dots, g(\nu_e)\big) \quad \text{is uniform on } \Fp^{t} \times \Fq^{e} .
\]
\end{lemma}

\begin{proof}
The tuple is the reduction of $g$ modulo $\prod_i (X - x_i) \cdot \prod_j m_j$, where $m_j$ is the minimal polynomial of $\nu_j$ over $\Fp$: indeed $\Fp[X]/(X - x_i) \cong \Fp$ by evaluation at $x_i$, and $\Fp[X]/(m_j) \cong \Fq$ by evaluation at $\nu_j$, with $\deg m_j = \extdeg$ ($\nu_j$ generates $\Fq$, since $\extdeg$ is prime and $\nu_j \notin \Fp$). The moduli are pairwise coprime: the $x_i$ are distinct, the $m_j$ are distinct irreducibles by non-conjugacy, and no $m_j$ has a root in $\Fp$. Their degree sum is $t + \extdeg\, e \leq 2^a$, so the reduction is surjective by Chinese remaindering. A surjective linear image of a uniform variable is uniform.
\end{proof}

\begin{proposition}[joint uniformity of the round-0 view]\label{prop:uniform}
Take $s_0 = 2$ and assume: the conditions of \Cref{thm:twopoint}; and that the $s_0 + s_1$ nodes (the commitment nodes $z_j^{2^{k_0}}$ and the out-of-domain points of $\hat{f}_1$) lie outside $\Fp$, are pairwise distinct and pairwise non-conjugate. Let $x_1, \dots, x_{t_0}$ be the distinct queried points (they lie in $\Fp \setminus \{0\}$, hence are automatically distinct from the nodes). Then the round-0 revealed values, namely the per-class evaluations $\hat{u}_s(\pow(x_t, n - k_0))$, the two commitment out-of-domain answers, the $2k_0$ fresh sumcheck coefficients, and the $s_1$ out-of-domain answers on $\hat{f}_1$, are jointly uniform, each on its natural space. In particular their joint law does not depend on the data.
\end{proposition}

\begin{proof}
Restrict every value, as an $\Fp$-linear form on the cells, to the union of the blocks. By \Cref{lem:blocks} the cross-terms vanish there, so each restricted form factors through the blockwise evaluations: the fiber value of class $s$ at $x_t$ reads the evaluation of the block univariate of $s$ at $x_t$, times the nonzero prefactor $x_t^{2^{n - k_0} - 2^a}$ (the blocks sit at the end of their classes), and the out-of-domain and sumcheck values apply the node matrix to the block evaluations at the nodes, by the computation of \Cref{sec:master}. The blockwise evaluation map is $\Fp$-surjective: per class by \Cref{lem:crt}, with the budget $t_0 + \extdeg(s_0 + s_1) \leq 2^a$, and jointly because the blocks are disjoint. The assembling map is $\Fq$-linear and surjective: the fiber coordinates pass through unchanged; the commitment out-of-domain and sumcheck values are the node matrix, of full row rank by \Cref{thm:twopoint}; and each $\hat{f}_1$ out-of-domain answer adds one row supported on the columns of its own node, where it is the nonzero vector $(\hat{eq}(\vec{\alpha}_0, s))_s$, so the extended matrix keeps full row rank. The values equal this composition applied to the uniform block cells, plus a contribution of the remaining cells (data and non-block mask cells): conditionally on the latter, they are a shifted surjective image of a uniform variable, hence uniform; so they are uniform unconditionally, with a law independent of the shift, hence of the data.
\end{proof}

\begin{lemma}[node genericity]\label{lem:nodeprob}
For uniform $z_1, z_2$ and uniform out-of-domain points of $\hat{f}_1$, the node hypotheses of \Cref{prop:uniform} fail with probability at most $(s_0 + s_1)\, p/q$ (a node in $\Fp$) plus $\binom{s_0 + s_1}{2}\, \extdeg\, 2^{k_0}/q$ (a collision or a Frobenius conjugation). For a commitment node the first count is exact: the set $\{z : z^{2^{k_0}} \in \Fp\}$ has size exactly $p$, since the $2^{k_0}$-th powers inside $\Fp^*$ form the subgroup of order $(p-1)/2^{k_0}$ (using $2^{k_0} \mid p - 1$) and each has $2^{k_0}$ preimages. For a pair of nodes: a Frobenius orbit has $\extdeg$ elements, each with at most $2^{k_0}$ preimages.
\end{lemma}

\subsection{The simulator}\label{sec:simulator}

Fix the statement (public parameters, public input, claimed value) and the verifier challenges. A \emph{data assignment} is a choice of contents for the data cells, consistent with the statement; the difference $\delta$ of two data assignments satisfies $\langle \delta, w \rangle = 0$ and vanishes on the public positions. Write $\mathrm{cells} = \mathrm{data} \oplus \mathrm{masks}$, and let $L := (L_{\mathrm{view}}, L_{\mathrm{fold}})$ be the $\Fp$-linear map sending the cells to the pair (round-0 view, $\hat{f}_1$), where $L_{\mathrm{view}}$ collects the values of \Cref{prop:uniform} and $L_{\mathrm{fold}}$ is the round-0 fold.

\begin{lemma}[reduction to round 0]\label{lem:reduction}
The whole transcript is a deterministic function of $L(\mathrm{cells})$, the statement, and the challenges.
\end{lemma}

\begin{proof}
The round-0 messages are the view together with deterministic completions (duplicate query answers, and the $X^0$ sumcheck coefficients, fixed by the running target). Every later message is computed by the honest prover from $\hat{f}_1$ and the challenges: the oracles $f_1, \dots, f_{M-1}$ are the encodings of the successive folds of $\hat{f}_1$, the later sumcheck messages, out-of-domain answers and openings are derived from these, and the final polynomial is the last fold.
\end{proof}

\begin{lemma}[exactness criterion]\label{lem:criterion}
Assume the conditions of \Cref{prop:uniform}. Let $W' := L_{\mathrm{fold}}\big(\ker(L_{\mathrm{view}}) \cap \mathrm{masks}\big)$. The transcripts produced with two data assignments differing by $\delta$ are identically distributed if and only if
\[
\varphi\big(L_{\mathrm{fold}}(\delta)\big) \;=\; \varphi\big(L_{\mathrm{fold}}(\varrho(L_{\mathrm{view}}(\delta)))\big) \qquad \text{for every } \varphi \in \mathrm{ann}(W') ,
\]
where $\varrho$ is any $\Fp$-linear section of $L_{\mathrm{view}}|_{\mathrm{masks}}$ (which is surjective by \Cref{prop:uniform}), and $\mathrm{ann}(W')$ is the space of $\Fp$-linear forms on the $\hat{f}_1$-space vanishing on $W'$.
\end{lemma}

\begin{proof}
By \Cref{lem:reduction} it suffices to compare the laws of $L(\mathrm{cells})$. With uniform masks, $L(\mathrm{cells})$ is uniform on a coset of the subspace $\mathcal{I} := L(\mathrm{masks})$, shifted by $L$ of the data; the two laws coincide if and only if $L(\delta) \in \mathcal{I}$. Since $L_{\mathrm{view}}|_{\mathrm{masks}}$ is surjective, $\mathcal{I} = \{(v, f) : f - L_{\mathrm{fold}}(\varrho(v)) \in W'\}$: for $(v,f) = L(m)$, the mask $m - \varrho(v)$ lies in $\ker(L_{\mathrm{view}}) \cap \mathrm{masks}$, and conversely. So $L(\delta) \in \mathcal{I}$ if and only if $L_{\mathrm{fold}}(\delta) - L_{\mathrm{fold}}(\varrho(L_{\mathrm{view}}(\delta))) \in W'$, which is the stated condition.
\end{proof}

\begin{lemma}[protocol directions]\label{lem:protocoldirs}
The following $\Fq$-forms on the $\hat{f}_1$-space (each counting as $\extdeg$ $\Fp$-forms) lie in $\mathrm{ann}(W')$ and satisfy the exactness criterion for \emph{every} $\delta$: (i) $f \mapsto \hat{f}(\pow(x_t, n - k_0))$ for each queried point $x_t$; (ii) $f \mapsto \hat{f}(\pow(z', n - k_0))$ for each out-of-domain point $z'$ of $\hat{f}_1$; (iii) $f \mapsto \sum_b \hat{f}(b)\, \hat{W}_0(\vec{\alpha}_0, b)$.
\end{lemma}

\begin{proof}
Each of these forms, composed with $L_{\mathrm{fold}}$, equals a fixed linear combination of the view coordinates, possibly plus a multiple of $\langle \cdot, w \rangle$, \emph{as an identity of linear forms on the cells}: (i) is the fold consistency $\hat{f}_1(\pow(x_t, \cdot)) = \fold(f_0, \vec{\alpha}_0)(x_t)$, where the right side is a public combination of the opened fiber of $x_t$; (ii) is the definition of the out-of-domain answer; (iii) is the terminal sumcheck identity $\sum_b \hat{f}_1(b)\, \hat{W}_0(\vec{\alpha}_0, b) = \hat{h}_{k_0}(\alpha_{k_0})$, whose right side is a public linear combination of the sumcheck values and of the initial target $\sigma_0 = y_1 + \gamma y_2 + \gamma^2 \langle \cdot, w \rangle$. Such a $\varphi$ vanishes on $\ker(L_{\mathrm{view}}) \cap \mathrm{masks}$ ($\langle \cdot, w \rangle$ vanishes on the masks), so $\varphi \in \mathrm{ann}(W')$; and both sides of its criterion equal the same combination applied to $L_{\mathrm{view}}(\delta)$, the term $\langle \cdot, w \rangle$ vanishing on data differences and on masks alike.
\end{proof}

\begin{condition}[no accidental pinning]\label{cond:pinning}
$\mathrm{ann}(W')$ is spanned by the forms of \Cref{lem:protocoldirs}.
\end{condition}

\begin{theorem}[perfect honest-verifier zero-knowledge, interactive]\label{thm:simulator}
Fix challenges satisfying the conditions of \Cref{prop:uniform} and \Cref{cond:pinning}. Let $D^*$ be any data assignment consistent with the statement (one is computable from the statement alone: fix the public positions, then meet the claimed value by adjusting one position carrying a nonzero weight). Then for every data assignment $D$, the honest transcript with data $D$ is distributed identically to the honest transcript with data $D^*$: the simulator runs the honest prover on $D^*$.
\end{theorem}

\begin{proof}
By \Cref{lem:criterion} it suffices to verify the criterion for $\delta := D - D^*$ and every $\varphi \in \mathrm{ann}(W')$; by \Cref{cond:pinning} it suffices for the forms of \Cref{lem:protocoldirs}, which satisfy it for every $\delta$.
\end{proof}

\begin{corollary}[statistical zero-knowledge]\label{cor:statzk}
For uniform challenges (the interactive setting), and for every data assignment $D$, the honest transcript with data $D$ is at statistical distance at most $\varepsilon$ from the output of the simulator (the honest prover run on $D^*$), with no prover-side check, where $\varepsilon := \Pr[\text{degenerate challenges}] \approx 2^{-122}$ (\Cref{cor:restart}).
\end{corollary}

\begin{proof}
The good set of challenge tuples (the conditions of \Cref{prop:uniform} and \Cref{cond:pinning}) is determined by the challenges and the statement, independently of the data. Conditioned on any good tuple, the two distributions coincide exactly (\Cref{thm:simulator}); on the bad set, bound the conditional distance by $1$. The bad set has the probability bounded in \Cref{cor:restart}.
\end{proof}

\begin{remark}[numerical verification]
We checked \Cref{cond:pinning} numerically at toy parameters ($p = 103$, $\extdeg = 3$, $n = 8$, $k_0 = 3$, $2^a = 16$, $t_0 = 2$, $s_0 = 2$, $s_1 = 1$), with the full machinery (fibers, out-of-domain answers, sumcheck coefficients including the cross-terms of a generic input weight $\hat{w}$, and the fold map), validated by passing the sumcheck recursion, the terminal identity, and \Cref{lem:blocks} as runtime assertions. Result, stable across seeds: $\dim \mathrm{ann}(W') = \extdeg\,(t_0 + s_1 + 1) = 12$, equal to the span of the protocol directions, which annihilate $W'$ exactly. Control experiment: with $\hat{w} = 0$, the dimension jumps to $15$, matching the prediction of \Cref{rem:identities} (the evaluations of $\hat{f}_1$ at the commitment nodes become pinned; the terminal direction already accounts for one of their combinations, leaving one new $\Fq$-direction).
\end{remark}

\subsection{Bounding the pinned directions}\label{sec:pinbound}

We now bound $\mathrm{ann}(W')$ from above, to control the probability that \Cref{cond:pinning} fails. Throughout this subsection, write $\lambda_s := \hat{eq}(\vec{\alpha}_0, s)$ for the coefficient of class $s$ in the round-0 fold (\Cref{sec:setting}).

\begin{definition}[\textsc{spread}]\label{def:spread}
The challenges $\vec{\alpha}_0$ satisfy \textsc{spread} if the $2^{k_0 - 1}$ fold multipliers $\lambda_{(1, s')}$, $s' \in \{0,1\}^{k_0 - 1}$, of the classes inside $R_{\mathsf{out}}$ span $\Fq$ as an $\Fp$-vector space.
\end{definition}

\begin{lemma}\label{lem:spread}
Let $\vec{\alpha}_0 = (\alpha_1, \dots, \alpha_{k_0})$ be independent uniform elements of $\Fq = \mathbb{F}_{p^{\extdeg}}$, with $\extdeg$ prime and $k_0 \geq \extdeg$. Then \textsc{spread} fails with probability at most
\[
\frac{1}{q} \;+\; \binom{k_0 - 1}{k_0 - \extdeg + 1} \Big(\frac{p}{q}\Big)^{k_0 - \extdeg + 1} ,
\]
which is at most $2/q$ for the practical parameters ($k_0 = 7$, $\extdeg = 5$: the second term is $20\, p^{-12} \leq p^{-5} = 1/q$).
\end{lemma}

\begin{proof}
Since $\hat{eq}(\alpha, 0) = 1 - \alpha$ and $\hat{eq}(\alpha, 1) = \alpha$, writing $A := \{i \geq 2 : s'_i = 1\}$ and expanding the factors $(1 - \alpha_i)$ gives
\[
\lambda_{(1, s')} \;=\; \alpha_1 \prod_{i \in A} \alpha_i \;+\; \text{($\pm 1$)-combination of the monomials } \alpha_1 \!\prod_{i \in S} \alpha_i \text{ with } S \supsetneq A .
\]
By downward induction on $|A|$, each monomial $\alpha_1 \prod_{i \in S} \alpha_i$, $S \subseteq \{2, \dots, k_0\}$, is in turn an integer combination of the family $\{\lambda_{(1, s')}\}_{s'}$: the two families have the same $\Fp$-span. If $\alpha_1 = 0$ (probability $1/q$) the span is $\{0\}$; otherwise multiplication by $\alpha_1$ preserves spanning, so it suffices to study $V_{k_0}$, where
\[
V_j \;:=\; \mathrm{span}_{\Fp}\Big\{ \prod_{i \in S} \alpha_i \;:\; S \subseteq \{2, \dots, j\} \Big\} , \qquad \text{so that } V_1 = \Fp \text{ and } V_j = V_{j-1} + \alpha_j V_{j-1} .
\]
We claim: if an $\Fp$-subspace $V$ satisfies $1 \in V$ and $V \neq \Fq$, then $T := \{x \in \Fq : xV \subseteq V\}$ equals $\Fp$. Indeed $T$ contains $\Fp$ and is closed under addition and multiplication; for $x \in T \setminus \{0\}$, multiplication by $x$ is injective on the finite set $T$, hence surjective, so $x^{-1} \in T$: $T$ is a subfield of $\Fq$ containing $\Fp$. Since $\extdeg$ is prime, $T = \Fp$ or $T = \Fq$, and the latter would give $V \supseteq \Fq \cdot 1 = \Fq$. So $T = \Fp$.

Consequently, whenever $V_{j-1} \neq \Fq$ and $\alpha_j \notin \Fp$, we have $\alpha_j V_{j-1} \not\subseteq V_{j-1}$, hence $\dim_{\Fp} V_j \geq \dim_{\Fp} V_{j-1} + 1$. Starting from $\dim_{\Fp} V_1 = 1$, if at least $\extdeg - 1$ of $\alpha_2, \dots, \alpha_{k_0}$ lie outside $\Fp$, then $V_{k_0} = \Fq$. Failure therefore requires at least $(k_0 - 1) - (\extdeg - 2) = k_0 - \extdeg + 1$ of them to lie in $\Fp$; each does so independently with probability $p/q$, and a union bound over the $\binom{k_0 - 1}{k_0 - \extdeg + 1}$ choices concludes.
\end{proof}

Throughout the rest of the subsection, fix challenges satisfying \textsc{spread}, $\alpha_1 \notin \{0, 1\}$, and the conditions of \Cref{prop:uniform}. Notation: by trace duality, an $\Fp$-form $\varphi$ on the $\hat{f}_1$-space $\Fq^{2^m}$ ($m := n - k_0$) has a unique representative $\phi \in \Fq^{2^m}$ with $\varphi(g) = \mathrm{tr}\,\langle \phi, g \rangle$, where $\langle \cdot, \cdot \rangle$ is the coordinatewise $\Fq$-pairing and $\mathrm{tr} = \mathrm{tr}_{\Fq/\Fp}$. Write $w^{(t)} \in \Fp^{2^m}$ for the weight vector of the evaluation at the queried point $x_t$, and $u^{(\nu)} \in \Fq^{2^m}$ for the weight vector of the evaluation at a node $\nu$ (the commitment nodes $\nu_1, \nu_2$ and the $s_1$ out-of-domain points of $\hat{f}_1$).

\begin{lemma}[confinement]\label{lem:confine}
Let $\varphi \in \mathrm{ann}(W')$ with representative $\phi$. Then, restricted to the block coordinates (the last $2^a$ positions),
\[
\phi\big|_{\mathrm{blk}} \;\in\; \mathrm{span}_{\Fq}\Big\{ w^{(t)}\big|_{\mathrm{blk}},\;\; \mathrm{slice}_r\big(u^{(\nu)}\big)\big|_{\mathrm{blk}} \;:\; t \leq t_0,\ \nu \text{ a node},\ r \leq \extdeg \Big\} ,
\]
where $\mathrm{slice}_r(u) := (\mathrm{tr}(\beta_r u_c))_c \in \Fp^{2^m}$ for a fixed $\Fp$-basis $(\beta_r)$ of $\Fq$.
\end{lemma}

\begin{proof}
Let $K \subseteq \Fp^{2^a}$ be the common kernel, inside the block coordinates, of the evaluation forms at the queried points and the nodes; by \Cref{lem:crt}, $\dim_{\Fp} K = 2^a - t_0 - \extdeg(s_0 + s_1)$. For any class $s$ and any $v \in K$, the perturbation placing $v$ on the block of class $s$ (and zero elsewhere) is invisible to the view: its fiber, out-of-domain, and $\hat{f}_1$-answer contributions vanish because all its point and node evaluations do, and its cross-term contributions vanish by \Cref{lem:blocks}. Its image under the fold is $\lambda_s v$ (placed at the block coordinates), so $\lambda_s v \in W'$ and $\mathrm{tr}(\lambda_s \langle \phi, v \rangle) = 0$. The multipliers $\lambda_s$ span $\Fq$ over $\Fp$: for the classes inside $R_{\mathsf{out}}$ this is \textsc{spread}, and it transfers to all classes since $\lambda_{(0, s')} = \frac{1 - \alpha_1}{\alpha_1}\, \lambda_{(1, s')}$ with $\alpha_1 \notin \{0, 1\}$. By nondegeneracy of the trace, $\langle \phi, v \rangle = 0$ for every $v \in K$. The vectors of an $\Fp$-basis of $K$ stay independent over $\Fq$, so the $\Fq$-orthogonal of $K$ in $\Fq^{2^a}$ has $\Fq$-dimension $t_0 + \extdeg(s_0 + s_1)$; it contains the stated span, whose generators are exactly an independent family of that size ($\Fp$-forms cutting $K$ out). The two spaces coincide, and $\phi|_{\mathrm{blk}}$ lies in them.
\end{proof}

\begin{lemma}[node channel]\label{lem:nodechannel}
Let $\varphi \in \mathrm{ann}(W')$ and write, by \Cref{lem:confine},
\[
\phi|_{\mathrm{blk}} \;=\; (\text{queried and } \hat{f}_1 \text{ parts}) \;+\; \sum_{j \leq s_0,\, r} b_{j,r}\, \mathrm{slice}_r\big(u^{(\nu_j)}\big)\big|_{\mathrm{blk}} , \qquad b_{j,r} \in \Fq .
\]
Define $B_{s, j} := \sum_r \mathrm{tr}(\lambda_s\, b_{j,r})\, \beta_r \in \Fq$. Then the vector $(B_{s,j})_{s,j}$ lies in the $\Fq$-row space of the node matrix of \Cref{sec:twopoint}.
\end{lemma}

\begin{proof}
For any assignment $V = (V_{s,j})$ in the $\Fq$-kernel of the node matrix, build the perturbation that realizes, on the block of each class $s$, the node values $(V_{s,1}, V_{s,2})$ and zero evaluations at the queried points and the $\hat{f}_1$ points: possible within the blocks by \Cref{lem:crt}. It is invisible to the view: fibers and $\hat{f}_1$ answers vanish by construction, cross-terms vanish on blocks, and the out-of-domain and sumcheck values vanish because $V$ kills their node combinations. Its fold image lies in $W'$, and pairing with $\varphi$ gives, using $\langle \mathrm{slice}_r(u^{(\nu_j)}), \delta_s \rangle = \mathrm{tr}(\beta_r V_{s,j})$ and $\mathrm{tr}(x\, \mathrm{tr}(y)) = \mathrm{tr}(x)\mathrm{tr}(y)$:
\[
0 \;=\; \varphi\Big(\sum_s \lambda_s \delta_s\Big) \;=\; \sum_{s, j, r} \mathrm{tr}(\lambda_s b_{j,r})\; \mathrm{tr}(\beta_r V_{s,j}) \;=\; \mathrm{tr}\Big( \sum_{s,j} B_{s,j} V_{s,j} \Big) .
\]
The kernel of the node matrix is an $\Fq$-subspace: replacing $V$ by $\theta V$ and varying $\theta \in \Fq$, nondegeneracy of the trace gives $\sum_{s,j} B_{s,j} V_{s,j} = 0$ for every kernel element $V$, that is, $B$ lies in the row space.
\end{proof}

The row space is spanned by the tensors $\{\omega_j, \tau_{\ell, j}\}$ (\Cref{thm:twopoint}), while $B$ ranges, as the $b_{j,r}$ vary, over the \emph{trace twists} of the extension rows: $B_{s,j} = T_{b_j}(\lambda_s)$ with $T_b(x) := \sum_r \mathrm{tr}(x\, b_r) \beta_r$, an arbitrary $\Fp$-linear map of $\Fq$. The untwisted case $T_b = (\text{multiplication by } \theta)$ corresponds to $\sum_j \theta_j \eta_j$, which does lie in the row space. The next two lemmas show that, generically, nothing else does, and that no untwisted combination is in turn pinned on its own: even the terminal combination carries a nonzero cross form, supplied, within the terminal protocol direction, by the input weight (\Cref{cond:cross2}).

\begin{lemma}[twist rigidity]\label{cond:twist}
Every $\Fp$-linear map of $\Fq$ is a linearized polynomial $T(x) = \sum_{r=0}^{\extdeg-1} t_r\, x^{p^r}$, and correspondingly $(T(\lambda_s))_s = \sum_r t_r\, \eta^{[r]}$, where $\eta^{[r]} := \bigotimes_i \big(1 - \alpha_i^{p^r},\; \alpha_i^{p^r}\big)$ is the extension row at the Frobenius-conjugated challenges (so $\eta^{[0]} = \eta$). Assume $\alpha_1 \notin \Fp$ and that, for each point $j$, the $(k_0 - 1) \times (\extdeg - 1)$ matrix
\[
M^{(j)} := \Big( \alpha_m^{p^r} - z_j^{2^{m-1}} \Big)_{m = 2..k_0,\; r = 1..\extdeg - 1}
\]
has rank $\extdeg - 1$. Then the only maps $T_{b_1}, T_{b_2}$ with $(T_{b_j}(\lambda_s))_{s,j}$ in the row space of the node matrix are the multiplications $T_{b_j} = \theta_j \cdot$, $\theta_j \in \Fq$.
\end{lemma}

\begin{proof}
The first claim: linearized polynomials are $\Fp$-linear, and they form a space of dimension $\extdeg^2$ over $\Fp$, equal to the dimension of all $\Fp$-linear maps, with no nonzero map represented twice (a nonzero linearized polynomial of degree $\leq p^{\extdeg - 1}$ has at most $p^{\extdeg-1} < q$ roots). For the tensor form: $\lambda_s^{p^r} = \prod_i \hat{eq}(\alpha_i, s_i)^{p^r} = \prod_i \hat{eq}(\alpha_i^{p^r}, s_i)$, since the Frobenius fixes the boolean values.

The row-space condition splits per point (the families $\{\omega_j, \tau_{\ell,j}\}$ are supported on their own blocks): for each $j$, $\sum_{r \geq 1} t_r\, \eta^{[r]} \in \mathrm{span}\{\omega, \tau_1, \dots, \tau_{k_0}\}$ at the point-$j$ slot values (the $r = 0$ term is absorbed by $\eta \in \mathrm{span}$, \Cref{prop:rank}). Fix $j$ and drop it. Expand all tensors in the per-slot basis $\{c_i, d\}$ (\cref{eq:cd}), with monomials indexed by the set $S$ of slots carrying $d$: writing $A^{[r]}_i := \alpha_i^{p^r} - z^{2^{i-1}}$, the coefficient of $\eta^{[r]}$ on $S$ is $\prod_{i \in S} A^{[r]}_i$ (\cref{eq:acd}); the coefficient of $\tau_m$ on $S$ is $\prod_{i \in S \setminus \{m\}} A^{[0]}_i$ if $m = \max S$, else $0$; the coefficient of $\omega$ is $1$ on $S = \emptyset$, else $0$. Suppose $\sum_{r \geq 1} y_r \eta^{[r]} = x_\omega \omega + \sum_m x_m \tau_m$. The singleton sets give $x_m = \sum_r y_r A^{[r]}_m$. The pairs $S = \{1, m\}$, $m \geq 2$, give $x_m A^{[0]}_1 = \sum_r y_r A^{[r]}_1 A^{[r]}_m$; substituting,
\[
\sum_{r \geq 1} \big( y_r\, D_r \big)\, A^{[r]}_m \;=\; 0 \qquad \text{for } m = 2, \dots, k_0 , \qquad D_r := A^{[r]}_1 - A^{[0]}_1 = \alpha_1^{p^r} - \alpha_1 .
\]
This says the vector $(y_r D_r)_r$ is in the kernel of $M^{(j)}$, which is trivial by the rank assumption. Since $\alpha_1 \notin \Fp$ and $\extdeg$ is prime, $D_r \neq 0$ for $1 \leq r \leq \extdeg - 1$ (the fixed field of the $r$-th Frobenius power inside $\Fq$ is $\Fp$). So $y_r = 0$ for all $r \geq 1$, and $T_{b_j}$ is the multiplication by its $t_0$-coefficient.
\end{proof}

\begin{corollary}\label{cor:twistprob}
For uniform challenges, the assumptions of \Cref{cond:twist} fail with probability at most
\[
\frac{p}{q} \;+\; 2 \binom{k_0 - 1}{\,k_0 - \extdeg + 1\,} \Big( \frac{p^{\extdeg - 2}}{q} \Big)^{k_0 - \extdeg + 1}
\qquad (\approx 2^{-124} \text{ for KoalaBear}) .
\]
\end{corollary}

\begin{proof}
$\Pr[\alpha_1 \in \Fp] = p/q$. For the rank, fix $j$ and process the rows of $M^{(j)}$ (one per $m$, each depending on its own $\alpha_m$) in order, tracking the span of the rows seen so far. Rank $< \extdeg - 1$ after $k_0 - 1$ rows requires at least $(k_0 - 1) - (\extdeg - 2) = k_0 - \extdeg + 1$ rows to fall into the span of their predecessors. A proper subspace of $\Fq^{\extdeg - 1}$ lies in a hyperplane $\{x : \sum_r v_r x_r = 0\}$ with $v \neq 0$, and the row of $\alpha_m$ falls in it only if $L_v(\alpha_m) = \sum_r v_r z^{2^{m-1}}$, where $L_v(x) := \sum_{r=1}^{\extdeg-1} v_r x^{p^r}$. Writing $L_v = \mathrm{Frob} \circ M_v$ with $M_v(x) = \sum_r v_r^{1/p} x^{p^{r-1}}$, a nonzero linearized polynomial of degree $\leq p^{\extdeg - 2}$, the solution set is empty or a coset of $\ker M_v$, of size at most $p^{\extdeg - 2}$. The rows are independent across $m$, so each fall costs at most $p^{\extdeg-2}/q$; a union bound over the choice of falling rows and over the two points gives the second term.
\end{proof}

\begin{lemma}[kernel surjectivity]\label{lem:kersurj}
Under the conditions of \Cref{prop:uniform}, every assignment $\delta_{\mathrm{out}}$ on the non-block mask cells extends to a perturbation in $\ker(L_{\mathrm{view}}) \cap \mathrm{masks}$.
\end{lemma}

\begin{proof}
The blocks must compensate the contributions of $\delta_{\mathrm{out}}$: its per-class fiber evaluations at the queried points, and its contributions to the $2k_0 + 2 + s_1$ out-of-domain, sumcheck, and $\hat{f}_1$-answer values, the latter through both node weights and cross weights. By \Cref{lem:crt}, the blocks realize any prescribed tuple of fiber and node evaluations, per class and jointly; and the assembled map from block node values to the $2k_0 + 2 + s_1$ values is surjective (full row rank, proof of \Cref{prop:uniform}). Choose block evaluations cancelling the fiber values classwise and the assembled values globally; blocks contribute no cross-terms (\Cref{lem:blocks}).
\end{proof}

\begin{lemma}[cross coupling]\label{cond:cross2}
Assume the conditions of \Cref{prop:uniform} and that $\hat{w}$ is nonzero at some data cell outside the blocks. Write $\eta_j = \sum_v c^{(j)}_v N_v$ for the unique representation of the extension rows in the value rows (\Cref{thm:twopoint}), and $F_\theta := \sum_v \big(\theta_1 c^{(1)}_v + \theta_2 c^{(2)}_v\big) C_v$ for the matching combination of the cross vectors, restricted to the non-block mask cells. Then, except with probability at most $2^{k_0 + 7}/q$ over the challenges, $F_\theta \neq 0$ for \emph{every} $\theta \neq 0$, including the terminal combination $\theta^{\mathrm{term}} := \big(P^{(1)}_{k_0+1},\, \gamma P^{(2)}_{k_0+1}\big)$, whose cross form is computed exactly in the proof; consequently, for every $\theta \neq 0$, the form $f \mapsto \theta_1 \hat{f}(\nu_1\text{-point}) + \theta_2 \hat{f}(\nu_2\text{-point})$ does not lie in $\mathrm{ann}(W')$.
\end{lemma}

\begin{proof}
First the consequence. For $\delta \in \ker(L_{\mathrm{view}}) \cap \mathrm{masks}$ with node values $V$ and non-block restriction $\delta_{\mathrm{out}}$, the vanishing of the value forms gives $N_v(V) = -C_v(\delta_{\mathrm{out}})$, hence
\[
\theta_1 \hat{f}(\nu_1) + \theta_2 \hat{f}(\nu_2) \Big|_{f = \mathrm{fold}(\delta)} \;=\; \sum_j \theta_j\, \eta_j(V) \;=\; \sum_{j, v} \theta_j c^{(j)}_v N_v(V) \;=\; - F_\theta(\delta_{\mathrm{out}}) ,
\]
and $\delta_{\mathrm{out}}$ is arbitrary by \Cref{lem:kersurj}: the form annihilates $W'$ if and only if $F_\theta = 0$.

Next, the terminal combination. Contrary to what the node-level identity $\eta = \rho_{k_0+1}$ may suggest, $F_{\theta^{\mathrm{term}}}$ does \emph{not} vanish: the identity lifts to the cells only together with the input-weight term of $\hat{W}_0$. Cell by cell, the terminal sumcheck identity $\sum_b \hat{f}_1(b)\, \hat{W}_0(\vec{\alpha}_0, b) = \hat{h}_{k_0}(\alpha_{k_0})$ decomposes as
\[
P^{(1)}_{k_0+1}\, \hat{f}_1(\nu_1\text{-point}) \;+\; \gamma\, P^{(2)}_{k_0+1}\, \hat{f}_1(\nu_2\text{-point}) \;+\; \gamma^2\, T_{\hat{w}} \;=\; \hat{h}_{k_0}(\alpha_{k_0}) , \qquad T_{\hat{w}} := \sum_{b} \hat{f}_1(b)\, \hat{w}(\vec{\alpha}_0, b) .
\]
As a linear form on the cells, $T_{\hat{w}}$ has coefficient $\lambda_s\, \hat{w}(\vec{\alpha}_0, c)$ at the cell of class $s$ and position $c$ (expanding $\hat{f}_0$ cell by cell, as in \Cref{sec:master}); at a block position $c$, $\hat{w}(\vec{\alpha}_0, c)$ is the multilinear extension, over the round-$0$ coordinates, of values that all vanish (\Cref{lem:blocks}), so, restricted to the masks, $T_{\hat{w}}$ is supported on the non-block mask cells. Evaluate the terminal identity on $\delta$: the right side is an explicit affine combination of the fresh sumcheck coefficients and of the initial target $\sigma_0 = y_1 + \gamma\, y_2 + \gamma^2 \langle \cdot, w \rangle$, and every term vanishes on $\delta$ (the values by definition of $\ker(L_{\mathrm{view}})$, the input term because $\hat{w}$ vanishes on the masks). Hence $\sum_j \theta^{\mathrm{term}}_j \eta_j(V) = -\gamma^2\, T_{\hat{w}}(\delta_{\mathrm{out}})$, and comparing with the display above, with $\delta_{\mathrm{out}}$ arbitrary:
\[
F_{\theta^{\mathrm{term}}} \;=\; \gamma^2\, T_{\hat{w}} \big|_{\text{non-block mask cells}} .
\]
(No contradiction with $\eta = \rho_{k_0+1}$: the \emph{pure} node combination $\sum_j \theta^{\mathrm{term}}_j \hat{f}(\nu_j\text{-point})$ is not in $\mathrm{ann}(W')$; the terminal \emph{protocol} direction is, because its input-weight component contributes exactly the missing cross form $\gamma^2 T_{\hat{w}}$. The control experiment $\hat{w} = 0$ of the numerical remarks is the degenerate case where $T_{\hat{w}}$, and every cross vector, vanishes.)

Injectivity. Since $\theta^{\mathrm{term}}_2 = \gamma P^{(2)}_{k_0+1} \neq 0$ under the standing conditions, the pair $\{(1,0),\, \theta^{\mathrm{term}}\}$ is a basis of $\Fq^2$: it suffices that $F_{(1,0)}$ and $F_{\theta^{\mathrm{term}}} = \gamma^2 T_{\hat{w}}$ be linearly independent over $\Fq$, which we witness on two cells. After clearing the denominators of the $c^{(1)}_v$ (Cramer over the value-row system, determinant nonzero under the standing conditions), each entry of $F_{(1,0)}$, and each $2 \times 2$ minor of the pair $(F_{(1,0)}, T_{\hat{w}})$, is a polynomial in the challenges of total degree at most $2^{k_0 + 7}$ (entries of the system have degree at most $2^{k_0}$ in each $z_j$, $k_0 + 1$ in the $\alpha$'s, $2$ in $\gamma$; a $16 \times 16$ determinant multiplies these by $16$; the entries of $T_{\hat{w}}$ add at most $2 k_0$). We exhibit a nonvanishing minor: specialize $\alpha_i := z_1^{2^{i-1}}$ for $i \geq 2$. Then $\rho^{(1)}_\ell = \omega^{(1)} + (\alpha_1 - z_1)\tau^{(1)}_1$ for every $\ell \geq 2$, and the system for $\eta_1$ solves explicitly: the block-1 equations for $\tau^{(1)}_m$, $m \geq 2$, give $\nu_{m,2} = -g(z_1^{2^{m-1}})\, \nu_{m,1}$; substituting into the block-2 equations for $\tau^{(2)}_m$ gives an upper-triangular system with diagonal $P^{(2)}_m (2z_2^{2^{m-1}} - 1)\, d_m \neq 0$, so $\nu_{m,1} = \nu_{m,2} = 0$ for all $m \geq 2$ by descent; the two remaining equations give
\[
\nu_{1,1} = \frac{\alpha_1 - z_1}{(2 z_1 - 1)\, d_1} , \qquad \nu_{1,2} = -\, g(z_2)\, \nu_{1,1} .
\]
The round-1 cross weights at a non-block mask cell $(1, v)$ are $[X^r]\big(X (1-X)\big)\, \gamma^2\, \hat{w}(0, v)$, that is $+\gamma^2 \hat{w}_0(v)$ for $r = 1$ and $-\gamma^2 \hat{w}_0(v)$ for $r = 2$, where $\hat{w}_0 := \hat{w}(0, \cdot)$. Hence, at the specialization,
\[
F_{(1,0)}(1, v) \;=\; \gamma^2\, \nu_{1,1} \big(1 + g(z_2)\big)\, \hat{w}_0(v) \;=\; \gamma^2\, \frac{(\alpha_1 - z_1)\; 2 z_2 (1 - z_2)}{(2z_1 - 1)\; d_1\; (2 z_2 - 1)}\; \hat{w}_0(v) ,
\]
and, at the mask cell $(1, v)$ with $v = (s', c)$ (class $(1, s')$, position $c$),
\[
T_{\hat{w}}(1, v) \;=\; \alpha_1 (1 - \alpha_1)\; e_{s'}\, A_c , \qquad e_{s'} := \prod_{i = 2}^{k_0} \hat{eq}\big(z_1^{2^{i-1}}, s'_i\big) , \qquad A_c := \sum_{t \in \{0,1\}^{k_0 - 1}} e_t\; \hat{w}_0(t, c) ,
\]
since, at the specialization, $\lambda_{(1, s')} = \alpha_1\, e_{s'}$ and $\hat{w}(\vec{\alpha}_0, c) = (1 - \alpha_1)\, A_c$ (the mask half contributes nothing to $\hat{w}(\vec{\alpha}_0, c)$). Let $(0, s^*, c^*)$ be the assumed data cell, outside the blocks, with $\hat{w}_0(s^*, c^*) \neq 0$, and pick any $s'' \neq s^*$ (possible: $k_0 \geq 2$). On the two cells $u := (1, (s^*, c^*))$ and $u'' := (1, (s'', c^*))$,
\[
F_{(1,0)}(u)\, T_{\hat{w}}(u'') - F_{(1,0)}(u'')\, T_{\hat{w}}(u) \;=\; \gamma^2\, \nu_{1,1} \big(1 + g(z_2)\big)\; \alpha_1 (1 - \alpha_1)\; A_{c^*} \Big[ \hat{w}_0(s^*, c^*)\, e_{s''} \;-\; \hat{w}_0(s'', c^*)\, e_{s^*} \Big] .
\]
Every factor is nonzero as a rational function of $(\alpha_1, z_1, z_2, \gamma)$: the first by the closed form above, the second trivially. For $A_{c^*}$ and the bracket, view them as multilinear polynomials in independent variables $\zeta_2, \dots, \zeta_{k_0}$: the evaluation at $\zeta := s^*$ gives $A_{c^*} = \hat{w}_0(s^*, c^*) \neq 0$, the evaluation at $\zeta := s''$ gives the bracket $= \hat{w}_0(s^*, c^*) \neq 0$; and a nonzero multilinear polynomial in $\zeta_2, \dots, \zeta_{k_0}$ stays nonzero under the substitution $\zeta_i := z_1^{2^{i-1}}$, since distinct multilinear monomials $\prod_{i \in S} \zeta_i$ land on distinct powers $z_1^{\sum_{i \in S} 2^{i-1}}$ (binary expansion). Schwartz--Zippel on the cleared minor concludes: the pair is independent, and $F_\theta \neq 0$ for every $\theta \neq 0$.
\end{proof}

\begin{remark}[slice level]\label{rem:slice}
\Cref{prop:pinbound} below applies \Cref{cond:cross2} to a single $\Fp$-form: what is used there is $\mathrm{tr} \circ F_\theta \neq 0$, slightly stronger than $F_\theta \neq 0$ (some \emph{entry} of $F_\theta$ must avoid the kernel of the trace). The determinant argument gives the $\Fq$-level statement; we leave an analytic bound for the slice-level event open. The numerical verification of \Cref{cond:pinning}, which computes $\mathrm{ann}(W')$ as an $\Fp$-space, operates directly at the slice level, and confirms it.
\end{remark}

\begin{proposition}\label{prop:pinbound}
Under \textsc{spread}, $\alpha_1 \notin \Fp$, the conditions of \Cref{prop:uniform}, and the assumptions of \Cref{cond:twist,cond:cross2}, the conclusion of \Cref{cond:pinning} holds: $\mathrm{ann}(W')$ is spanned by the protocol directions.
\end{proposition}

\begin{proof}
Let $\varphi \in \mathrm{ann}(W')$. By \Cref{lem:confine,lem:nodechannel,cond:twist}, the commitment-node part of $\phi|_{\mathrm{blk}}$ is the trace dual of $\sum_j \theta_j \eta_j$ for scalars $\theta_j$, and the queried-point and $\hat{f}_1$-node parts are $\Fq$-combinations of protocol directions outright. Subtract the latter: the residue $\varphi'$ lies in $\mathrm{ann}(W')$, with block representative the dual of $\sum_j \theta_j \eta_j$; and \Cref{lem:confine}'s orthogonality argument (whose perturbations of the $R_{\mathsf{out}}$ classes range over the entire class, so the conclusion also constrains the non-block coordinates) leaves, beyond it, no component except along the input-weight direction $\gamma^2\, \hat{w}(\vec{\alpha}_0, \cdot)$, the one the terminal protocol direction itself carries: $\varphi'$ is the trace dual of $\sum_j \theta_j\, u^{(\nu_j)} + e\, \gamma^2\, \hat{w}(\vec{\alpha}_0, \cdot)$ for some $e \in \Fq$. Evaluating on $W'$, with the identity $F_{\theta^{\mathrm{term}}} = \gamma^2\, T_{\hat{w}}$ from the proof of \Cref{cond:cross2}: for every $\delta \in \ker(L_{\mathrm{view}}) \cap \mathrm{masks}$,
\[
0 \;=\; \varphi'\big(\mathrm{fold}(\delta)\big) \;=\; \mathrm{tr}\Big( \sum_j \theta_j\, \eta_j(V) \;+\; e\, \gamma^2\, T_{\hat{w}}(\delta_{\mathrm{out}}) \Big) \;=\; -\,\mathrm{tr}\Big( F_{\theta - e\, \theta^{\mathrm{term}}}(\delta_{\mathrm{out}}) \Big) ,
\]
with $\delta_{\mathrm{out}}$ arbitrary (\Cref{lem:kersurj}). By \Cref{cond:cross2}, applied at the slice level (\Cref{rem:slice}), $\theta = e\, \theta^{\mathrm{term}}$: the residue is $e$ times the terminal direction, and $\varphi$ is a combination of protocol directions.
\end{proof}

\begin{corollary}[degenerate challenges]\label{cor:restart}
A challenge tuple is outside the good set only if one of the following fails: \textsc{spread} (at most $2/q$, \Cref{lem:spread}); the conditions of \Cref{prop:uniform} (at most $(4 \cdot 2^{k_0} + 2k_0)/q$ for \Cref{thm:twopoint}, plus $(s_0 + s_1)\, p/q + \binom{s_0 + s_1}{2} \extdeg\, 2^{k_0}/q$ for the node conditions, \Cref{lem:nodeprob}); the twist-rigidity assumptions (at most $p/q$ plus a negligible term, \Cref{cor:twistprob}); the cross coupling (at most $2^{k_0+7}/q$, \Cref{cond:cross2}, at the $\Fq$ level: \Cref{rem:slice}). Total, for KoalaBear with $s_0 = 2$:
\[
\Pr[\mathrm{degenerate\ challenges}] \;\leq\; \big(s_0 + s_1 + 1\big)\, \frac{p}{q} \;+\; O\big(2^{k_0 + 7}\big)\frac{1}{q} \;\approx\; 2^{-122} ,
\]
dominated by the events that a sampled point or node lands in the base field.
\end{corollary}

\begin{remark}
The numerical verification of \Cref{cond:pinning} reported above doubles as an end-to-end test of \Cref{prop:pinbound}: at the toy parameters, $\dim \mathrm{ann}(W')$ equals the protocol count for generic $\hat{w}$ across seeds, and the control $\hat{w} = 0$, which violates the hypothesis of \Cref{cond:cross2}, produces exactly one extra pinned direction, as the analysis predicts. That the generic count is exactly the protocol count, and not more, also confirms $F_{\theta^{\mathrm{term}}} \neq 0$: were the pure terminal node combination pinned, it would contribute an extra $\Fq$-direction, independent of the protocol span. The explicit solution in the proof of \Cref{cond:cross2} was additionally verified at production shape ($k_0 = 7$, sixteen rows over the $512$ node values, toy field): at the specialization, the solved representation of $\eta_1$ matches the stated $\nu_{1,1}$, $\nu_{1,2}$, $\mu_1$, $\mu_2$ exactly, with all other coefficients zero, and the closed-form factor agrees.
\end{remark}

\section{Logup-GKR sumcheck}\label{sec:zk-logup-gkr}

TODO

\section{AIR sumcheck}\label{sec:zk-air-sumcheck}

TODO

\section{Column evaluations}\label{sec:zk-column-evals}

In the two protocols above (\Cref{sec:zk-logup-gkr,sec:zk-air-sumcheck}), the prover sends a few evaluations, at random points, of the columns, interpreted as multilinear polynomials storing data we want to hide; for some columns $c$ it also sends an evaluation of $\mathsf{shift}(c)$.

Strategy:
\begin{itemize}
    \item For each table $T$ and each column $c$ that will be evaluated $t$ times and can take values in $\Omega$ (in the sense that the other columns can be filled so that the AIR constraints hold), add $r = t \cdot \log_2(q / |\Omega|) + 2 \lambda$ rows whose values in column $c$ are chosen independently and uniformly at random in $\Omega$; the rest of the trace is adapted to satisfy the constraints.
    \item By the conjecture below, these structured random masks protect the column evaluations: statistical zero-knowledge at $\lambda$ bits.
    \item Conjecture: let $P$ be a multilinear polynomial in $n$ variables, $r = t \cdot \log_2(q / |\Omega|) + 2 \lambda$ of whose coefficients (in the Lagrange basis) are sampled independently and uniformly at random, the other $2^n - r$ set to zero. Evaluate it at $t$ random points, each sampled independently and uniformly in $\Fq^n$, some of the evaluations possibly taken on $\mathsf{shift}(P)$ instead of $P$. Then the $t$ evaluations are at statistical distance at most $2^{-\lambda}$ from the uniform distribution over $\Fq^t$.
\end{itemize}

\textit{No protocol changes.}

\bibliographystyle{IEEEtran}
\bibliography{bibliography}

\end{document}
